% CPSC 4900 Senior Project Proposal -- Olivia Ye, Fall 2026
% Structure follows the example CPSC proposals: background, related work,
% project description, evaluation, timeline, deliverables, references.

\documentclass[11pt]{article}
\usepackage[margin=0.85in]{geometry}
\usepackage{enumitem}
\usepackage{booktabs}
\usepackage{array}
\usepackage{listings}
\usepackage[compact]{titlesec}
\usepackage{multicol}
\usepackage[hidelinks]{hyperref}

\titlespacing{\section}{0pt}{4pt}{1pt}
\setlist{itemsep=1pt, topsep=2pt, parsep=0pt}
\setlength{\parskip}{0pt}
\linespread{0.97}

\lstset{
  basicstyle=\ttfamily\footnotesize,
  frame=single, framesep=3pt, xleftmargin=6pt,
  columns=fullflexible, keepspaces=true,
  aboveskip=4pt, belowskip=4pt
}

\title{\vspace{-1.8em}Risk Unwind: Typed Exception Handling for\\ Portfolio Risk Management}
\author{Olivia Ye \\ \small Advised by Professor Stephen Slade}
\date{CPSC 4900 \\ Fall 2026}

\begin{document}
\maketitle
\vspace{-2.2em}

\section{Introduction}

Programming languages have a mature, well-defined way of dealing with things going wrong:
raise an exception, match it to a handler, resolve it, and if nothing can handle it, stop and
escalate. Portfolio risk management has the detection half of that machinery and not the
dispatch half. Commercial systems notice limit breaches reliably and route them to a person;
nothing decides, from the \emph{kind} of problem, what the right response is.

\begin{lstlisting}[language=Python]
try:
    portfolio.hold()
except (ConcentrationBreach, VaRBreach) as breaches:
    propose(solve(cost, subject_to=every_limit))  # net, then solve once
except DataFault as f:
    quarantine(f.feed); mark_stale(f.metrics)     # the numbers aren't real
except ModelIntegrityFault as f:
    switch_to(simpler_measure); flag(f.model)     # a broken model is worse
except LiquidityFault as b:
    raise Uncatchable(b) from b
finally:
    verify(portfolio)                             # every path, including failure
\end{lstlisting}

Four branches, four genuinely different responses, and only one of them trades. That is the
claim: not that breaches can be fixed automatically, but that \emph{what kind of problem it
is} determines whether fixing is even the right move.

\section{Background}

A risk limit is a promise about a portfolio: no more than a fixed share in any one name, no
more than a certain loss on a bad day, no position larger than the market can absorb. Markets
move, so these promises break constantly without anyone doing anything wrong: an assumption
the system relied on is no longer true, which is exactly what an exception is. The
correspondence is structural --- the \emph{type} is what kind of promise broke,
\emph{uncaught} is escalation with the full trace, \emph{finally} verifies the repair --- but
four properties of financial risk make it harder than a direct port:

\begin{enumerate}[leftmargin=*]
\item \textbf{Breaches cascade.} A stock sits inside a sector inside a region, so one
  oversized position trips three limits: three alarms, one event.
\item \textbf{Breaches can genuinely offset.} Two positions may each breach alone while
  partially hedging each other, leaving the combined book inside its bound. Subadditivity is
  one of the four axioms of a coherent risk measure~\cite{artzner1999}: risk is formally not
  additive. Cascades are structural, offsets mathematical, and they need different handling.
\item \textbf{Some breaches cannot be fixed.} If a position is oversized and its market
  has dried up, no trade solves the problem.
\item \textbf{Some breaches must not be fixed.} If the feed has died or the model has
  broken down, the breach may not be real, and trading on it is worse than doing nothing.
  The correct response is to stop --- which no optimiser will ever choose, because stopping
  is not in its vocabulary.
\end{enumerate}

\section{The Hierarchy, and Why It Is the Contribution}

Risk limits today live in flat lists: one rule per limit, one response per rule, and the
response is nearly always ``tell someone.'' Exceptions form a family tree instead, and that
structure is what does the work.

\begin{lstlisting}[basicstyle=\ttfamily\scriptsize]
RiskException
|-- LimitBreach                      |-- DataFault
|   |-- ConcentrationBreach          |-- ModelIntegrityFault
|   |   |-- SingleNameBreach         \-- LiquidityFault
|   |   \-- SectorBreach
|   \-- VaRBreach
\end{lstlisting}

A handler is written at whatever level the response is actually the same --- one
\texttt{ConcentrationBreach} handler covers every sector and instrument --- and the type raised
selects it, so \emph{what to do} is decided by the taxonomy rather than by the check that
detected the problem.

\begin{center}
\begin{tabular}{@{}>{\raggedright\arraybackslash}p{2.9cm}>{\raggedright\arraybackslash}p{8.8cm}@{}}
\toprule
\textbf{Type} & \textbf{Response} \\
\midrule
\texttt{Concentration}/ \texttt{VaRBreach} & Net the breaches, then compute one joint set
of trades clearing all of them at least cost. \textbf{Trades.} \\
\texttt{DataFault} & Quarantine the feed, mark dependent metrics stale, stop.
\textbf{Trades nothing} --- any trade would act on fiction. \\
\texttt{ModelIntegrity\-Fault} & Fall back to a simpler, robust measure; flag the model
for validation. \textbf{Trades nothing} until the fallback is in place. \\
\texttt{LiquidityFault} & No trade fixes this. Propagate outward and escalate with the
causal chain. \textbf{Cannot trade.} \\
\bottomrule
\end{tabular}
\end{center}

An optimiser knows exactly one move: trade. Hand it a dead price feed and it will confidently
rebalance against numbers that are not real. This is an exception system rather than a solver
with monitors bolted to the front precisely because three of the four branches resolve by
\emph{declining to act}, and something must decide which branch applies.

\textbf{Proposes, does not execute.} Every handler emits a proposal --- the trade, its cost,
the breaches it clears, the reasoning --- for a person to approve, since Rule 15c3-5 requires
risk controls to sit under a broker-dealer's direct and exclusive control~\cite{sec15c35}.
Resolution is not autonomous; what changes is the message to that person, from \emph{``four
limits broke, go work out what to do''} to \emph{``these four alarms are one event, here is
the smallest trade that clears it, and here is what it will not fix.''}

\section{Related Work}

Detection is a solved commercial problem, and any proposal claiming otherwise will not survive
a reader who works in the industry. BlackRock's Aladdin checks compliance at order creation,
at execution and after the trade, and runs a workflow in which violations are assigned,
commented on and aged~\cite{aladdin}; pre-trade controls are legally required for
broker-dealers with market access~\cite{sec15c35}. But at the moment of a breach these
systems hand it to a human: that is detection and \emph{reporting}, not exception
\emph{handling} --- no handler computes a resolution, and the \texttt{except} block is an
inbox.

Automated remediation is not unclaimed either. US patents 7,636,686 and 7,930,239~\cite{patent686}
resolve a restriction violation by executing a stored rule that substitutes another security or
cash for the restricted holding, and 7,599,869 describes pre-trade compliance checking with an
override path and an aged exception queue. But that is per-restriction substitution at order
time --- one rule, one violation, one response --- not reasoning about several live breaches
together, and not deciding from the type not to trade at all.

The computer science lineage supplies the machinery: Goodenough on what a handler must guarantee
and the difference between resuming and unwinding~\cite{goodenough1975}, Russell et al.\ on
exception patterns in business processes~\cite{russell2006}, and MAPE-K~\cite{kephart2003} for
the monitor--detect--plan--act shape. Basel already defines seven operational risk event
types~\cite{basel}, so this project is not first to classify risk events --- it argues that
categories built for operational losses do not fit limit breaches.

That leaves four claims, and only four. \textbf{(1)} Dispatch is by type, so three of four
classes resolve by refusing to trade --- which no optimiser, and none of the patents above,
can do. \textbf{(2)} Breaches are reasoned about jointly, so cascaded alarms collapse to one
action and offsetting ones to none. \textbf{(3)} When no trade can fix the problem, that
arrives as a mathematical result, not a case someone remembered to write. \textbf{(4)} The
evaluation is reproducible, on injected ground truth rather than a private book.

\section{System Design}

\textbf{Main goal:} build and evaluate a typed exception handling system for a single
equity portfolio, which raises typed exceptions from replayed market data, dispatches on
type, computes a minimal joint repair where one is warranted, declines to trade where it
is not, and escalates what it cannot fix.

\textbf{Testable hypothesis:} typed dispatch plus joint netting resolves a given set of
risk events with strictly less trading than a per-breach handler, and avoids the trades a
type-free optimiser makes on corrupted inputs, with no loss in detection accuracy.

Six stages, built in this order because each is only testable once the one before it
exists. Nothing here is a neural network; no part of the project involves one.

\begin{itemize}[leftmargin=*]
\item \textbf{S0 --- Replay.} Historical data on a simulated clock, so a month of trading runs
  in seconds, plus a harness that \emph{injects} faults deliberately: break a limit, cut a
  feed, make two breaches coincide. Because the faults are deliberate the correct answer is
  known, which is the only reason anything downstream can be graded.
\item \textbf{S1 --- State.} On each tick, update value, exposures and risk measures ---
  bookkeeping, but every later stage reads from it.
\item \textbf{S2 --- Typed monitors.} Independent monitors, each raising a \emph{typed}
  exception rather than a generic alert and reporting how far out of bounds, not merely that
  something broke. A monitor never decides what to do --- that separation is Goodenough's
  central idea, and it lets remediation policy change without touching a detector.
\item \textbf{S3 --- Net, decode, link.} Recompute risk with offsetting positions combined and
  discard breaches that disappear; match the survivors against per-cause fingerprints,
  tolerating a missing or spurious alarm, following practice in
  telecommunications~\cite{jakobson1993}; record causal chains. Netting must precede decoding,
  or the decoder is fed alarms that were never real.
\item \textbf{S4 --- Repair.} Reached only for \texttt{LimitBreach} types; if a data or
  model fault is live, dispatch never arrives here, which is the point of having types.
  Repair minimises trading cost subject to every active limit holding afterwards and no trade
  exceeding what the market can absorb. A worked \emph{cascade}: in a \$10M portfolio capped
  at 5\% per holding and 25\% per sector, a chip stock at \$780k is \$280k over while
  semiconductors at 27\% are \$200k over --- but trimming the stock drags the sector to
  24.2\%, so \$280k of trading clears both, not \$480k. Every limit enters as a constraint,
  not only the broken ones, which stops a repair from creating a fresh breach; and when they
  cannot jointly be met, the solver proves it.
\item \textbf{S5 --- Verify and escalate.} Re-run the monitors on the proposed portfolio;
  anything still breached, plus anything infeasible, goes to a person with the full trace.
\end{itemize}

A solver is used rather than a learned model because a breach is arithmetic and the repair a
constrained optimisation: the answer provably satisfies every constraint, is auditable as model
risk governance requires, and proves infeasibility when no repair exists.

\section{Evaluation}

Eight constructed scenarios, each with a known correct answer, layered on replayed data from
real stress periods so market behaviour is genuine even though the injected faults are not.

\begin{enumerate}[leftmargin=*]
\item \textbf{A corrupted feed alongside a genuine breach --- the headline test.} The correct
  action is to trade \emph{nothing} and freeze; a system without types rebalances against
  fabricated numbers. Only typed dispatch wins this one, and the failure mode is catastrophic
  rather than merely wasteful.
\item A single isolated breach (detected at all?); two unrelated breaches at once (two
  independent repairs); a containment cascade, stock inside sector inside region (one fix, not
  three); two genuinely offsetting breaches (one fix, or none); a liquidity constraint making
  another repair infeasible (is the link found?); a breach with no feasible repair (does it
  escalate rather than fail quietly?); a model integrity fault during live breaches (does it
  fall back?).
\end{enumerate}

Three baselines: a plain limit checker with no correlation logic; a handler repairing every
breach separately, which should visibly over-trade on the cascade; and a bare optimiser with no
type system, which treats every alarm as a trade signal and fails hardest on scenario 1.
Metrics, in order: \textbf{dispatch correctness} (right type raised, right handler fired), then
detection and false-alarm rates, correct-cause identification, repair cost against the optimum,
latency, and false escalation rate.

\textbf{Scope, stated deliberately:} one long/short equity portfolio; simulated time, not
production streaming infrastructure; equities only, no non-linear payoffs; a solver, not a
learned policy; proposals reviewed by a person, never autonomous execution; roughly six leaf
types, where the test of the tree is that every split is defensible, not complete.

\section{Timeline}

{\small
\begin{itemize}[leftmargin=*]
\item \textbf{Sep.\ 22} --- Proposal submitted to advisor for review and approval.
  \textbf{Sep.\ 29} --- literature survey complete: a systematic search of IEEE Xplore, the
  ACM Digital Library and SSRN, and the three patents read in full, before any novelty claim
  is written. Dataset selected. Establish how often genuine offsets survive in practice,
  since many systems net at the point of \emph{measurement} --- this sets how much of S3 is
  needed.
\item \textbf{Oct.\ 6} --- S0 replay and fault injection working; S1 reads from it.
  \textbf{Oct.\ 13} --- hierarchy fixed, S2 typed monitors, scenarios 2--3 pass end to end;
  midterm progress report submitted.
\item \textbf{Oct.\ 20} --- Dispatch layer; scenarios 1 and 8 pass (declines to trade on a
  corrupted feed; falls back on a model fault). \textbf{Oct.\ 27} --- S3 netting, decoding
  and linking; scenario 4 shows a measurable trade-count difference against the per-breach
  baseline.
\item \textbf{Nov.\ 10} --- S4 solver with infeasibility reporting; scenarios 5--7 pass.
  \textbf{Nov.\ 17} --- S5 verification, proposal objects and escalation traces; full
  pipeline runs unattended. \textbf{Nov.\ 24} --- all three baselines run, every metric
  collected.
\item \textbf{Dec.\ 1} --- results compiled, poster and presentation drafted.
  \textbf{Dec.\ 8} --- draft to advisor. \textbf{Dec.\ 14} --- final report submitted.
\end{itemize}}

\section{Deliverables}

{\small
\begin{enumerate}[leftmargin=*]
\item Source code for the full S0--S5 pipeline as a public GitHub repository, with the exception
  hierarchy and dispatch table documented as a standalone artifact.
\item Fault injection harness, the eight evaluation scenarios and the three baselines, all
  reproducible from a single command, plus a results notebook holding every metric table and
  figure behind the final report.
\item Poster, final written report, ten-minute video presentation, and both documents added to
  the CPSC 4900 website.
\end{enumerate}}

\scriptsize
\begin{multicols}{2}
\begin{thebibliography}{9}
\setlength{\itemsep}{0pt}
\bibitem{goodenough1975} J. B. Goodenough. Exception handling: issues and a proposed notation.
  \emph{Communications of the ACM}, 18(12):683--696, 1975.
\bibitem{russell2006} N. Russell, W. M. P. van der Aalst, and A. H. M. ter Hofstede. Workflow
  exception patterns. In \emph{CAiSE 2006}, LNCS 4001:288--302, 2006.
\bibitem{artzner1999} P. Artzner, F. Delbaen, J.-M. Eber, and D. Heath. Coherent measures of
  risk. \emph{Mathematical Finance}, 9(3):203--228, 1999.
\bibitem{kephart2003} J. O. Kephart and D. M. Chess. The vision of autonomic computing.
  \emph{Computer}, 36(1):41--50, 2003.
\bibitem{sec15c35} U.S. SEC. Risk Management Controls for Brokers or Dealers with Market
  Access, 17 C.F.R. \S 240.15c3-5. Rel.\ 34-63241, adopted Nov.\ 3, 2010; 75 Fed.\ Reg.\ 69791.
\bibitem{patent686} Automated actions based on restrictions. U.S. Patents 7,636,686 and
  7,930,239, CheckFree Corp. See also U.S. Patent 7,599,869.
\bibitem{basel} Basel Committee on Banking Supervision. Operational risk event types. % TODO: pin exact document -- Basel II Annex 9, carried into the
  % consolidated Basel Framework (OPE).
\bibitem{jakobson1993} G. Jakobson and M. Weissman. Alarm correlation. \emph{IEEE Network},
  7(6):52--59, 1993. % TODO: volume/issue/pages not yet verified.
\bibitem{aladdin} BlackRock. Aladdin compliance and risk management platform.
  % TODO: cite the specific product page or documentation, with an access date.
\end{thebibliography}
\end{multicols}

\end{document}
